
\chapter{Les matrices}
\label{chap-matrices}

\section{Définition}

\begin{definition}[Matrice]
Soient $m \in \N^*$ et $n \in \N^*$. Une
\emph{matrice} à $m$ lignes et $n$ colonnes est un
tableau rectangulaire de $m \times n$ réels : il y a
$m$ lignes, et chaque ligne contient $n$ éléments.
\[
    A =
    \begin{pmatrix}
        a_{11} & a_{12} & \cdots & a_{1n} \\
        a_{21} & a_{22} & \cdots & a_{2n} \\
        \vdots & \vdots & \ddots & \vdots \\
        a_{m1} & a_{m2} & \cdots & a_{mn}
    \end{pmatrix},
    \qquad a_{ij} \in \R.
\]
Le réel $a_{ij}$ est le \emph{coefficient} de la
ligne $i$ et de la colonne $j$. On note aussi
$A = (a_{ij})_{1 \le i \le m,\, 1 \le j \le n}$.
L'ensemble des matrices à $m$ lignes et $n$ colonnes
est noté $\R^{m \times n}$, et le couple $(m, n)$
est la \emph{taille} de $A$.
\end{definition}

\begin{remark}
On donne toujours le nombre de lignes \emph{avant}
le nombre de colonnes, aussi bien pour la taille
$m \times n$ que pour l'indice $a_{ij}$. Une matrice
$2 \times 3$ et une matrice $3 \times 2$ ne sont pas
de la même taille.
\end{remark}

\begin{definition}[Lignes et colonnes]
Pour $A \in \R^{m \times n}$ :
\begin{itemize}
    \item la \emph{ligne} $i$ de $A$ est
          $(a_{i1}, a_{i2}, \dots, a_{in})$ : elle a
          $n$ éléments ;
    \item la \emph{colonne} $j$ de $A$ est
          $(a_{1j}, a_{2j}, \dots, a_{mj})^\top$ :
          elle a $m$ éléments.
\end{itemize}
\end{definition}

\begin{definition}[Égalité]
Deux matrices $A$ et $B$ sont égales si elles ont
\textbf{la même taille} et si $a_{ij} = b_{ij}$ pour
tous $i$ et $j$.
\end{definition}

\section{Matrices particulières}

\begin{definition}[Matrice carrée]
Une matrice est \emph{carrée} si elle a autant de
lignes que de colonnes ($m = n$). Ses coefficients
$a_{11}, a_{22}, \dots, a_{nn}$ forment sa
\emph{diagonale}.
\end{definition}

\begin{definition}[Matrice nulle]
La matrice nulle de $\R^{m \times n}$, notée $0$
(ou $0_{m,n}$), a tous ses coefficients égaux à $0$.
\end{definition}

\begin{definition}[Matrice identité]
La matrice identité $I_n \in \R^{n \times n}$ est la
matrice carrée qui a des $1$ sur la diagonale et des
$0$ partout ailleurs :
\[
    I_n =
    \begin{pmatrix}
        1 & 0 & \cdots & 0 \\
        0 & 1 & \cdots & 0 \\
        \vdots & \vdots & \ddots & \vdots \\
        0 & 0 & \cdots & 1
    \end{pmatrix},
    \qquad
    (I_n)_{ij} =
    \begin{cases}
        1 & \text{si } i = j, \\
        0 & \text{sinon.}
    \end{cases}
\]
\end{definition}

\begin{remark}
Un vecteur $\vect{x} \in \R^n$ peut être vu comme une
matrice à $n$ lignes et une seule colonne, c'est-à-dire
un élément de $\R^{n \times 1}$. C'est cohérent avec la
convention de ce carnet : les vecteurs sont des
vecteurs colonnes.
\end{remark}

Dans toute la suite, $A$ et $B$ sont deux matrices
de $\R^{m \times n}$ et $k \in \R$ est un réel.

\section{Opérations}

\subsection{Addition et soustraction}

\begin{definition}
Pour deux matrices \textbf{de même taille}, on
additionne (ou soustrait) coefficient par
coefficient :
\[
    (A + B)_{ij} = a_{ij} + b_{ij},
    \qquad
    (A - B)_{ij} = a_{ij} - b_{ij}.
\]
Le résultat est encore une matrice de
$\R^{m \times n}$.
\end{definition}

\begin{importantbox}
La somme de deux matrices de tailles différentes
n'est pas définie, même si elles ont le même nombre
total de coefficients (par exemple une $2 \times 3$
et une $3 \times 2$).
\end{importantbox}

\subsection{Multiplication par un réel}

\begin{definition}
\[
    (kA)_{ij} = k\, a_{ij}.
\]
Chaque coefficient est multiplié par $k$. Comme pour
les vecteurs, la division par un réel $k \neq 0$ est
la multiplication par $\frac{1}{k}$.
\end{definition}

\subsection{Transposée}

\begin{definition}[Transposée]
La \emph{transposée} de $A \in \R^{m \times n}$ est
la matrice $A^\top \in \R^{n \times m}$ définie par
\[
    (A^\top)_{ij} = a_{ji}.
\]
Les lignes de $A$ deviennent les colonnes de
$A^\top$, et inversement.
\end{definition}

\begin{importantbox}
La transposée échange les dimensions : si $A$ a
$m$ lignes et $n$ colonnes, $A^\top$ a $n$ lignes et
$m$ colonnes. Elle ne garde la même taille que pour
une matrice carrée.
\end{importantbox}

\begin{proposition}
Pour $A, B \in \R^{m \times n}$ et $k \in \R$ :
\begin{enumerate}
    \item $(A^\top)^\top = A$ ;
    \item $(A + B)^\top = A^\top + B^\top$ ;
    \item $(kA)^\top = k\, A^\top$.
\end{enumerate}
\end{proposition}

\begin{definition}[Matrice symétrique]
Une matrice carrée $A$ est \emph{symétrique} si
$A^\top = A$, c'est-à-dire si $a_{ij} = a_{ji}$ pour
tous $i$ et $j$. Par exemple, $I_n$ est symétrique.
\end{definition}

\subsection{Produit de Hadamard}

\begin{definition}[Produit terme à terme]
Le \emph{produit de Hadamard} de deux matrices
\textbf{de même taille} est la matrice
$A \odot B \in \R^{m \times n}$ définie par
\[
    (A \odot B)_{ij} = a_{ij}\, b_{ij}.
\]
\end{definition}

\begin{importantbox}
Ce n'est \emph{pas} le produit matriciel de la
\cref{def-produit-matriciel}. Attention à
NumPy~\cite{numpy} : l'opérateur \texttt{*} y calcule
$A \odot B$, et c'est \texttt{@} (ou
\texttt{np.dot}) qui calcule $AB$.
\end{importantbox}

\subsection{Produit matriciel}

\begin{definition}[Produit matriciel]
\label{def-produit-matriciel}
Soient $A \in \R^{m \times n}$ et
$B \in \R^{n \times p}$. Le \emph{produit}
$AB \in \R^{m \times p}$ est défini par
\[
    (AB)_{ij} = \sum_{k=1}^{n} a_{ik}\, b_{kj}
    = a_{i1} b_{1j} + a_{i2} b_{2j} + \dots + a_{in} b_{nj}.
\]
Le coefficient $(AB)_{ij}$ est le produit scalaire
de la ligne $i$ de $A$ et de la colonne $j$ de $B$.
\end{definition}

\begin{importantbox}
Le produit $AB$ n'existe que si le nombre de
\textbf{colonnes} de $A$ est égal au nombre de
\textbf{lignes} de $B$ :
\[
    \underbrace{(m \times n)}_{A}
    \cdot
    \underbrace{(n \times p)}_{B}
    = \underbrace{(m \times p)}_{AB}.
\]
Les deux dimensions « intérieures » doivent être
égales, et les deux dimensions « extérieures » donnent
la taille du résultat.
\end{importantbox}

\begin{proposition}
Pour des matrices dont les tailles rendent les
produits possibles, et pour $k \in \R$ :
\begin{enumerate}
    \item associativité : $(AB)C = A(BC)$ ;
    \item distributivité : $A(B + C) = AB + AC$ et
          $(A + B)C = AC + BC$ ;
    \item $(kA)B = A(kB) = k(AB)$ ;
    \item élément neutre : pour
          $A \in \R^{m \times n}$,
          $I_m A = A I_n = A$ ;
    \item transposée d'un produit :
          $(AB)^\top = B^\top A^\top$.
\end{enumerate}
\end{proposition}

\begin{importantbox}
Le produit matriciel n'est \textbf{pas commutatif} :
en général $AB \neq BA$. Même quand les deux
produits existent, ils peuvent ne pas avoir la même
taille, et même quand $A$ et $B$ sont carrées, ils
sont en général différents.
\end{importantbox}

\subsection{Produit d'une matrice par un vecteur}

\begin{definition}[Produit matrice-vecteur]
\label{def-matrice-vecteur}
Soient $A \in \R^{m \times n}$ et $\vect{x} \in \R^n$.
En voyant $\vect{x}$ comme une matrice de
$\R^{n \times 1}$, le produit $A\vect{x}$ est un
vecteur de $\R^m$ :
\[
    (A\vect{x})_i = \sum_{k=1}^{n} a_{ik}\, x_k,
    \qquad i \in \{1, \dots, m\}.
\]
La composante $i$ de $A\vect{x}$ est le produit
scalaire de la ligne $i$ de $A$ avec $\vect{x}$.
\end{definition}

\begin{remark}
Le produit scalaire est lui-même un produit
matriciel : pour $\vect{u}, \vect{v} \in \R^n$,
$\vect{u}^\top \in \R^{1 \times n}$ et
$\vect{v} \in \R^{n \times 1}$, donc
$\vect{u}^\top \vect{v} \in \R^{1 \times 1}$ est la
matrice qui contient le réel
$\vect{u} \cdot \vect{v}$.
\end{remark}

\subsection{Trace}

\begin{definition}[Trace]
La \emph{trace} d'une matrice carrée
$A \in \R^{n \times n}$ est la somme de ses
coefficients diagonaux :
\[
    \operatorname{tr}(A) = \sum_{i=1}^{n} a_{ii}.
\]
Par exemple, $\operatorname{tr}(I_n) = n$.
\end{definition}

\section{Exemples calculés à la main}

Tous les exemples utilisent les trois matrices
\[
    A = \begin{pmatrix} 1 & 2 \\ 3 & 4 \end{pmatrix},
    \qquad
    B = \begin{pmatrix} 0 & 1 \\ 1 & 0 \end{pmatrix},
    \qquad
    C = \begin{pmatrix} 1 & 2 & 3 \\ 4 & 5 & 6 \end{pmatrix}.
\]
$A$ et $B$ sont des matrices $2 \times 2$, et $C$ est
une matrice $2 \times 3$.

\begin{example}[Addition, soustraction et multiplication par un réel]
\[
    A + B = \begin{pmatrix} 1 & 3 \\ 4 & 4 \end{pmatrix},
    \qquad
    A - B = \begin{pmatrix} 1 & 1 \\ 2 & 4 \end{pmatrix},
    \qquad
    2A = \begin{pmatrix} 2 & 4 \\ 6 & 8 \end{pmatrix}.
\]
En revanche, $A + C$ n'est pas défini : $A$ est
$2 \times 2$ et $C$ est $2 \times 3$.
\end{example}

\begin{example}[Transposées]
\[
    A^\top = \begin{pmatrix} 1 & 3 \\ 2 & 4 \end{pmatrix},
    \qquad
    C^\top = \begin{pmatrix} 1 & 4 \\ 2 & 5 \\ 3 & 6 \end{pmatrix}.
\]
$C^\top$ est une matrice $3 \times 2$ : la première
ligne de $C$ est devenue sa première colonne.
$B^\top = B$, donc $B$ est symétrique.
\end{example}

\begin{example}[Produit de Hadamard et trace]
\[
    A \odot B =
    \begin{pmatrix} 1 \times 0 & 2 \times 1 \\ 3 \times 1 & 4 \times 0 \end{pmatrix}
    = \begin{pmatrix} 0 & 2 \\ 3 & 0 \end{pmatrix},
    \qquad
    \operatorname{tr}(A) = 1 + 4 = 5,
    \qquad
    \operatorname{tr}(B) = 0.
\]
\end{example}

\begin{example}[Le produit n'est pas commutatif]
\[
    AB =
    \begin{pmatrix}
        1 \times 0 + 2 \times 1 & 1 \times 1 + 2 \times 0 \\
        3 \times 0 + 4 \times 1 & 3 \times 1 + 4 \times 0
    \end{pmatrix}
    = \begin{pmatrix} 2 & 1 \\ 4 & 3 \end{pmatrix},
    \qquad
    BA = \begin{pmatrix} 3 & 4 \\ 1 & 2 \end{pmatrix}.
\]
On a bien $AB \neq BA$. Multiplier par $B$
\emph{à droite} échange les colonnes de $A$, et
multiplier par $B$ \emph{à gauche} échange ses
lignes. $B$ est une \emph{matrice de permutation} :
on la retrouvera au \cref{chap-systemes}.
\end{example}

\begin{example}[Tailles des produits]
Pour chaque produit, on compare les dimensions
intérieures :
\begin{center}
\begin{tabular}{@{}lccc@{}}
    \toprule
    Produit & Tailles & Existe ? & Taille du résultat \\
    \midrule
    $AC$ & $(2 \times 2)(2 \times 3)$ & oui & $2 \times 3$ \\
    $CA$ & $(2 \times 3)(2 \times 2)$ & non : $3 \neq 2$ & --- \\
    $C^\top A$ & $(3 \times 2)(2 \times 2)$ & oui & $3 \times 2$ \\
    $C C^\top$ & $(2 \times 3)(3 \times 2)$ & oui & $2 \times 2$ \\
    $C^\top C$ & $(3 \times 2)(2 \times 3)$ & oui & $3 \times 3$ \\
    \bottomrule
\end{tabular}
\end{center}
Par exemple,
$C C^\top = \begin{pmatrix} 14 & 32 \\ 32 & 77 \end{pmatrix}$ :
le coefficient $(1, 1)$ vaut
$1^2 + 2^2 + 3^2 = 14$, c'est le carré de la norme de
la première ligne de $C$.
\end{example}

\begin{example}[Transposée d'un produit]
\[
    AC =
    \begin{pmatrix}
        1 + 8 & 2 + 10 & 3 + 12 \\
        3 + 16 & 6 + 20 & 9 + 24
    \end{pmatrix}
    = \begin{pmatrix} 9 & 12 & 15 \\ 19 & 26 & 33 \end{pmatrix},
    \qquad
    (AC)^\top = \begin{pmatrix} 9 & 19 \\ 12 & 26 \\ 15 & 33 \end{pmatrix}.
\]
On vérifie le premier coefficient de $C^\top A^\top$ :
la ligne $1$ de $C^\top$ est $(1, 4)$, la colonne $1$
de $A^\top$ est $(1, 2)^\top$, et
$1 \times 1 + 4 \times 2 = 9$. On retrouve bien
$(AC)^\top = C^\top A^\top$. L'ordre doit s'inverser :
$A^\top C^\top$ n'existe même pas, car
$(2 \times 2)(3 \times 2)$ n'a pas les bonnes
dimensions intérieures.
\end{example}

\begin{example}[Produit d'une matrice par un vecteur]
\[
    A \begin{pmatrix} 1 \\ 1 \end{pmatrix}
    = \begin{pmatrix} 1 + 2 \\ 3 + 4 \end{pmatrix}
    = \begin{pmatrix} 3 \\ 7 \end{pmatrix}.
\]
Multiplier par le vecteur $(1, 1)^\top$ fait la somme
de chaque ligne.
\end{example}

\section{Utilité dans le projet}

Comme dans NNFS~\cite{nnfs}, une couche de neurones
reçoit un \emph{lot} de $N$
exemples. On range ces exemples dans une matrice
$X \in \R^{N \times n}$ : une ligne par exemple, une
colonne par entrée. Les poids de la couche forment une
matrice $W \in \R^{n \times p}$ : une ligne par entrée
et une colonne par neurone. La sortie de la couche,
avant la fonction d'activation, est
\[
    Z = X W + \vect{b},
    \qquad
    \underbrace{(N \times n)}_{X}
    \cdot
    \underbrace{(n \times p)}_{W}
    = \underbrace{(N \times p)}_{Z}.
\]
Le coefficient $z_{ij}$ est le produit scalaire de
l'exemple $i$ et des poids du neurone $j$ : on retrouve
le calcul du neurone seul
(\cref{sec-vecteurs-utilite}), fait
d'un coup pour tous les neurones et tous les exemples.
L'ajout du biais $\vect{b} \in \R^p$ à chaque ligne
de $XW$ est le \emph{broadcasting}.

La transposée servira à la rétropropagation : le
gradient par rapport aux poids s'écrit avec $X^\top$.
Le produit matriciel est l'opération la plus appelée
du projet, et c'est pour cette raison que sa vitesse
a été mesurée.

\section{Implémentation}
\label{sec-matrices-implementation}

\subsection{Stockage à plat}

Une matrice $m \times n$ est rangée dans un
\textbf{seul} tableau de $m \times n$ réels, ligne
après ligne (\emph{rangement par lignes}, comme en C++
et dans NumPy). Avec les indices du C++, qui
commencent à $0$, l'élément de la ligne $i$ et de la
colonne $j$ se trouve à la position
\[
    i \times n + j.
\]
Pour la matrice $C$ ci-dessus, le tableau est
$(1, 2, 3, 4, 5, 6)$ et l'élément $c_{23} = 6$
(indices C++ $i = 1$, $j = 2$) est à la position
$1 \times 3 + 2 = 5$.

\begin{remark}
Un \cpp|vector<vector<double>>| ferait une
allocation par ligne, avec des lignes dispersées en
mémoire, et rien ne garantirait qu'elles aient toutes
la même longueur~\cite{cppreference}. Le tableau à
plat fait une seule allocation, et des éléments
voisins dans une ligne sont voisins en mémoire.
\end{remark}

\subsection{Ordre des boucles du produit}

La formule $(AB)_{ij} = \sum_k a_{ik} b_{kj}$ se code
avec trois boucles imbriquées sur $i$, $j$ et $k$.
L'ordre de ces boucles ne change pas le résultat, mais
il change beaucoup le temps de calcul :
\begin{itemize}
    \item dans l'ordre $i, j, k$, la boucle intérieure
          fait varier $k$ et lit $b_{kj}$ en sautant
          de $p$ cases à chaque tour, c'est-à-dire
          une colonne de $B$ ;
    \item dans l'ordre $i, k, j$, la boucle intérieure
          fait varier $j$ : elle lit une ligne de $B$
          et écrit une ligne du résultat, dans des
          cases voisines en mémoire.
\end{itemize}
Le processeur charge la mémoire par blocs de 64 octets
(8 réels). Dans le second ordre, chaque bloc chargé
sert 8 fois, et le compilateur peut traiter plusieurs
réels par instruction.

\begin{importantbox}
Mesure du 16 novembre, produit de deux matrices
$500 \times 500$ (environ $2 \times 500^3$, soit
250 millions d'opérations), compilé en \texttt{-O2} :
environ \qty{59}{\ms} dans l'ordre $i, j, k$ et \qty{21}{\ms} dans
l'ordre $i, k, j$, soit près de 3 fois plus rapide.
\end{importantbox}

\begin{codebox}
Le produit matriciel de MiniNumpy, dans l'ordre
$i, k, j$. La fonction privée \cpp|idx(i, j)|
calcule la position $i \times n + j$ dans le tableau
à plat.
\end{codebox}

\lstinputlisting[
language=C++,
caption=Produit matriciel (\texttt{src/Matrix.cpp}),
linerange=debutProduit-finProduit,
numbers=left,
breaklines=true
]{../../src/Matrix.cpp}

\newpage
